\documentclass[11pt]{article}
\usepackage[a4paper,margin=1in]{geometry}
\usepackage{amsmath,amssymb,amsthm,booktabs,enumitem,microtype,tikz}
\usepackage[colorlinks=true,linkcolor=blue!55!black,citecolor=blue!55!black,urlcolor=blue!55!black]{hyperref}
\newtheorem{theorem}{Theorem}[section]
\newtheorem{proposition}[theorem]{Proposition}
\newtheorem{lemma}[theorem]{Lemma}
\newtheorem{corollary}[theorem]{Corollary}
\theoremstyle{definition}
\newtheorem{definition}[theorem]{Definition}
\newtheorem{remark}[theorem]{Remark}
\newcommand{\bip}{\operatorname{bip}}
\newcommand{\supp}{\operatorname{supp}}
\newcommand{\mc}{\operatorname{MaxCut}}

\title{Shortest-geodesic supports in triangle-free maximum cuts:\\
an infinite Hall obstruction and the first minimal footprints}
\author{Alper Ferudun\\\texttt{alper@mercurycodelab.com}}
\date{July 2026}

\begin{document}
\maketitle

\begin{abstract}
Fix a cut of a triangle-free graph, write \(B\) for its crossing edges and
\(M\) for its monochromatic edges, and associate to each \(xy\in M\) the
\(B\)-edges lying on shortest \(x\)-\(y\) paths. A natural local approach
to triangle-free bipartization asks whether these supports satisfy Hall
expansion. We show that they need not, even when the cut is the unique
maximum cut and every bad edge has a unique shortest geodesic.

For every \(t\ge1\), we construct a triangle-free graph \(G_t\) on
\(7t+3\) vertices with a unique maximum cut and \(\bip(G_t)=t^2\).
Its \(t^2\) bad edges all have cut-distance four, but their support union
has only \(2t+2\) edges. Thus Hall expansion fails for \(t\ge3\), with an
unbounded ratio. An explicit packing of \(t^2\) edge-disjoint \(7\)-cycles
certifies the maximum cut and its uniqueness. We also derive the finite
footprint forced by a minimal Hall violation. Exact enumeration gives no
footprint with at most eight atoms, one footprint with nine atoms, and
three footprint isomorphism classes with ten atoms. The graph \(G_3\)
realizes the unique nine-atom obstruction. The family theorem is
combinatorial; the finite classification is computer-assisted and has two
independent exact verifiers. These results do not resolve the Erdos
\(n^2/25\) conjecture.
\end{abstract}

\section{Introduction}

For a finite graph \(G\), let
\[
 \bip(G)=|E(G)|-\mc(G)
\]
be its edge-deletion distance to bipartiteness. Erdos, Faudree, Pach and
Spencer asked whether every triangle-free \(N\)-vertex graph satisfies
\[
 \bip(G)\le N^2/25. \tag{1}\label{eq:erdos}
\]
The balanced blow-up of \(C_5\) shows that the constant would be sharp.
See \cite{EFPS88,EGS92,BCL21} for the problem and the strongest general
bounds. The conjectured value is known for \(N=5n\), \(1\le n\le40\),
by an exact certificate \cite{Ferudun26}.

We study a local packing principle suggested by a fixed maximum cut.
Let \(B\) be the crossing subgraph and \(M\) the edges left monochromatic.
For \(xy\in M\), the endpoints lie on the same shore. If \(B\) is
connected, then \(d_B(x,y)\) is even; triangle-freeness excludes distance
two. The first case is therefore \(d_B(x,y)=4\). One may try to assign a
distinct shortest-path edge to every such bad edge by Hall's theorem.

Our main theorem says that this principle fails by an arbitrarily large
factor even under unique maximum-cut and unique-geodesic hypotheses. The
use of shortest odd cycles in max-cut estimates has precedents, notably
\cite{CKL17}, but the full shortest-geodesic support incidence system
considered here appears to be different.

\section{Supports and minimal violations}

Fix a cut \(V(G)=V_0\mathbin{\dot\cup}V_1\), and set
\[
 B=\{xy\in E(G):x\in V_0,\ y\in V_1\},\qquad M=E(G)\setminus B.
\]
Assume that \(B\) is connected.

\begin{definition}
For \(e=xy\in M\), define
\[
 \supp_B(e)=\{f\in B:f\text{ lies on a shortest }x\text{-}y
 \text{ path in }B\}.
\]
For \(\mathcal S\subseteq M\), let
\(\supp_B(\mathcal S)=\bigcup_{e\in\mathcal S}\supp_B(e)\).
The length-five bad edges have \emph{support Hall expansion} when
\[
 |\supp_B(\mathcal S)|\ge|\mathcal S| \tag{2}\label{eq:hall}
\]
for every \(\mathcal S\subseteq\{xy\in M:d_B(x,y)=4\}\).
\end{definition}

Condition \eqref{eq:hall} is Hall's condition for the incidence graph
between bad edges and the crossing edges in their shortest supports.

\begin{proposition}[Support-size gap]\label{prop:gap}
If \(B\) is bipartite and \(d_B(x,y)=4\), then
\[
 |\supp_B(xy)|=4\quad\text{or}\quad|\supp_B(xy)|\ge6.
\]
\end{proposition}

\begin{proof}
If the shortest path is unique, its four edges form the support. Otherwise
choose two distinct shortest paths \(P,Q\). Their nonempty symmetric
difference is an even-degree subgraph and contains a cycle. Since \(B\) is
bipartite, this cycle has at least four edges. Hence
\[
 |E(P)\cup E(Q)|
 =\frac{4+4+|E(P)\triangle E(Q)|}{2}\ge6.
\]
\end{proof}

\begin{lemma}[Minimal footprint]\label{lem:minimal}
Let \(G\) be triangle-free, and let \(\mathcal S\) be inclusion-minimal
with
\[
 |\supp_B(\mathcal S)|<|\mathcal S|=m.
\]
Writing \(F=\supp_B(\mathcal S)\), we have:
\begin{enumerate}[label=\textup{(\roman*)}]
\item \(F\) is connected and bipartite, with \(|E(F)|=m-1\);
\item every edge of \(F\) belongs to at least two atom supports;
\item every \(xy\in\mathcal S\) has distance four in \(F\), and its
support is the union of all length-four \(F\)-geodesics from \(x\) to \(y\);
\item the graph \((V(F),\mathcal S)\) is triangle-free.
\end{enumerate}
\end{lemma}

\begin{proof}
For every \(e\in\mathcal S\), minimality gives
\[
 |\supp_B(\mathcal S\setminus\{e\})|\ge m-1.
\]
This set is contained in \(F\), while \(|E(F)|<m\). Hence
\(|E(F)|=m-1\) and removing any atom leaves the whole support union.
This proves (ii).

Every atom support is connected. If \(F\) had several components, the
atoms partitioned among them, and one component would itself contain more
atoms than support edges, contradicting minimality. Thus \(F\) is
connected. It is bipartite because \(F\subseteq B\).

A shortest \(B\)-path for an atom lies in \(F\), so its \(F\)-distance is
at most four. The distance is even, and distance two would make a triangle
with the bad edge. Thus the distance is four. Every length-four path in
\(F\) is also a shortest path in \(B\), proving (iii). Finally,
\(\mathcal S\subseteq E(G)\), proving (iv).
\end{proof}

\section{An infinite Hall obstruction}

Let \(L,A,B,C,D,E,R\) be seven disjoint independent sets of size \(t\),
and let \(u,w,v\) be three further vertices. Define \(G_t\) by adding
all edges between consecutive classes in
\[
 L-A-B-C-D-E-R,
\]
all edges of \(L-R\), and the thin channel
\[
 L-u-w-v-R,
\]
where \(L-u\) and \(v-R\) are complete blocks. Then
\[
 |V(G_t)|=7t+3,\qquad |E(G_t)|=7t^2+2t+2.
\]
The graph is triangle-free: neither channel has a chord, and no vertex
outside \(L\cup R\) is adjacent to both terminal classes.

\begin{theorem}\label{thm:family}
For every \(t\ge1\), \(G_t\) has a unique maximum cut up to global
complementation, and
\[
 \bip(G_t)=t^2,\qquad \mc(G_t)=6t^2+2t+2.
\]
The bad edges are exactly \(L-R\).
\end{theorem}

\begin{proof}
Index each large class by \(\mathbb Z_t\). For every
\((i,j)\in\mathbb Z_t^2\), form the \(7\)-cycle
\[
 L_iA_jB_{i+j}C_iD_jE_{i+j}R_jL_i. \tag{3}
\]
These \(t^2\) cycles are edge-disjoint. On the six wide-channel blocks,
their edge labels are
\[
 (i,j),\ (j,i+j),\ (i+j,i),\ (i,j),\ (j,i+j),\ (i+j,j),
\]
and each map is a bijection of \(\mathbb Z_t^2\). Their closing
\(L_iR_j\) edges are also distinct.

Every cut leaves a monochromatic edge on each packed odd cycle, so it has
at least \(t^2\) bad edges. Conversely, put \(L,B,D,R,w\) on one shore
and \(A,C,E,u,v\) on the other. Exactly the \(L-R\) edges are bad.

For uniqueness, a maximum cut has exactly one bad edge in each packed
cycle and no bad edge outside the packing. The outside edges are the thin
channel. Since all \(L_i u,uw,wv,vR_j\) cross, all vertices in \(L\cup R\)
lie on one shore and \(u,v\) on the other. Thus all \(t^2\) edges of
\(L-R\) are bad, and every remaining edge crosses. This determines the
cut up to complementation.
\end{proof}

\begin{corollary}\label{cor:hall}
In the unique maximum cut of \(G_t\), every bad edge \(L_iR_j\) has the
unique shortest \(B\)-geodesic
\[
 L_i-u-w-v-R_j,
\]
and
\[
 |\supp_B(M)|=2t+2.
\]
Thus \eqref{eq:hall} fails for \(t\ge3\), and
\[
 \frac{|M|}{|\supp_B(M)|}=\frac{t^2}{2t+2}\longrightarrow\infty.
\]
\end{corollary}

\begin{proof}
After deleting \(L-R\), the thin channel has length four and the wide
channel has length six. The thin internal vertices are singletons, so the
length-four path is unique. The support union consists of the \(t\) edges
\(L-u\), the two edges \(uw,wv\), and the \(t\) edges \(v-R\).
\end{proof}

The maximum cut is unique and therefore automatically minimizes every
secondary tie-break among maximum cuts. In particular,
\[
 \Gamma(B,M)=\sum_{xy\in M}(d_B(x,y)+1)^2=25t^2.
\]
The graphs \(G_t\) are not counterexamples to \eqref{eq:erdos}, since
\(t^2<(7t+3)^2/25\).

\section{Exact classification of the first footprints}

\begin{definition}
An \(m\)-footprint obstruction is a pair \((F,\mathcal A)\) such that:
\begin{enumerate}[label=\textup{(\alph*)}]
\item \(F\) is connected and bipartite with \(m-1\) edges;
\item \(\mathcal A\) is a set of \(m\) distinct distance-four pairs in \(F\);
\item the graph \((V(F),\mathcal A)\) is triangle-free;
\item the geodesic supports \(P_F(a)\) cover \(E(F)\), and every edge of
\(F\) lies in at least two supports.
\end{enumerate}
\end{definition}

Lemma~\ref{lem:minimal} sends every graph-level minimal Hall violation to
an \(m\)-footprint obstruction.

\begin{theorem}[Computer-assisted]\label{thm:classification}
No \(m\)-footprint obstruction exists for \(m\le8\). For \(m=9\), there
is exactly one footprint isomorphism class and its atom set is forced:
\[
 \{l_1,l_2,l_3\}-u-w-v-\{r_1,r_2,r_3\},
 \qquad
 \mathcal A=\{l_ir_j:1\le i,j\le3\}. \tag{4}
\]
For \(m=10\), there are exactly three footprint isomorphism classes.
\end{theorem}

The cases \(m\le6\) have a short direct proof. A minimal violation has
\(m-1\) support edges. Every atom support has at least four edges, so
\(m\ge5\). At \(m=5\), the footprint is a four-edge path and determines
only one distance-four pair. At \(m=6\), Proposition~\ref{prop:gap}
forces each support to be one of the at most \(\binom54=5\) four-edge
subsets of a five-edge footprint; a path edge set determines its endpoint
pair.

For \(7\le m\le10\), the proof is exhaustive. Using \texttt{geng} from
nauty/Traces \cite{McKayPiperno14}, we generated one representative of
every connected bipartite graph with \(m-1\) edges and between \(5\) and
\(m\) vertices. Breadth-first search listed every distance-four pair and
its full geodesic support. We then enumerated every \(m\)-element set of
pairs and tested conditions (c) and (d). Table~\ref{tab:counts} gives the
exact output. A second, independently written analysis program reproduced
the table.

\begin{table}[ht]
\centering
\begin{tabular}{c@{\quad}r@{\quad}r@{\quad}r}
\toprule
\(m\) & distance-filtered footprints & valid footprint classes
& admissible atom subsets\\
\midrule
6&0&0&0\\
7&0&0&0\\
8&0&0&0\\
9&1&1&1\\
10&9&3&56\\
\bottomrule
\end{tabular}
\caption{Exact footprint enumeration. The final column counts subsets on
canonical footprint representatives, before quotienting those subsets by
footprint automorphisms.}
\label{tab:counts}
\end{table}

The graph6 label in (4) is \texttt{H???FaM}. The ten-atom footprint labels
are \texttt{I????B\_fo}, \texttt{I????Bobo}, and
\texttt{I???E?wh\_}. Consequently, every graph-level minimal violation has
at least nine atoms, and the nine-atom footprint is uniquely determined.
The graph \(G_3\) realizes it. More generally, choosing three vertices
from each of \(L\) and \(R\) in any \(G_t\), \(t\ge3\), yields the same
minimal subsystem.

\section{Interpretation and reproducibility}

The wide channel in \(G_t\) certifies maximum-cut optimality by packing
odd cycles, while the thin channel forces all bad edges to share the same
two central support edges. Thus the resources certifying cut optimality
and the resources offered by shortest supports can live on different
channels.

This rules out bare shortest-support Hall expansion, not stronger banked
inequalities that also use vertex slack, longer geodesics, or cut-boundary
capacity. It also gives no stronger numerical upper bound on \(\bip(G)\).

The ancillary archive contains two family verifiers, two footprint enumerators, a command-level README, and a SHA-256 manifest. The family theorem is independent of computation.
The classification uses integer arithmetic throughout and relies on
nauty/Traces only for complete unlabeled graph generation.

\begin{thebibliography}{9}
\bibitem{BCL21}
J.~Balogh, F.~C.~Clemen and B.~Lidicky,
\emph{Max cuts in triangle-free graphs},
Combinatorica 43 (2023), 1059--1079;
\href{https://arxiv.org/abs/2103.14179}{arXiv:2103.14179}.

\bibitem{CKL17}
S.~M.~Cioaba, J.~H.~Koolen and W.~Li,
\emph{Max-cut and extendability of matchings in distance-regular graphs},
European J. Combin. 62 (2017), 232--244;
\href{https://arxiv.org/abs/1507.06254}{arXiv:1507.06254}.

\bibitem{EFPS88}
P.~Erdos, R.~Faudree, J.~Pach and J.~H.~Spencer,
\emph{How to make a graph bipartite},
J. Combin. Theory Ser. B 45 (1988), 86--98,
\href{https://doi.org/10.1016/0095-8956(88)90057-3}{doi:10.1016/0095-8956(88)90057-3}.

\bibitem{EGS92}
P.~Erdos, E.~Gyori and M.~Simonovits,
\emph{How many edges should be deleted to make a triangle-free graph
bipartite?},
in \emph{Sets, Graphs and Numbers}, Colloq. Math. Soc. Janos Bolyai 60,
North-Holland, 1992, 239--263.

\bibitem{Ferudun26}
A.~Ferudun,
\emph{The Erdos \(n^2/25\) max-cut conjecture for small multiples of five,
via a per-root-MaxCut envelope and blow-up integrality},
\href{https://arxiv.org/abs/2606.28041}{arXiv:2606.28041}, 2026.

\bibitem{McKayPiperno14}
B.~D.~McKay and A.~Piperno,
\emph{Practical graph isomorphism, II},
J. Symbolic Comput. 60 (2014), 94--112,
\href{https://doi.org/10.1016/j.jsc.2013.09.003}{doi:10.1016/j.jsc.2013.09.003}.
\end{thebibliography}
\end{document}



